\documentclass[12pt]{article}
\usepackage[utf8]{inputenc}
\usepackage[T1]{fontenc}
\usepackage[margin=1in]{geometry}
\usepackage{setspace}
\usepackage{xcolor}
\usepackage[colorlinks=true, urlcolor=blue!60!black, linkcolor=blue!60!black]{hyperref}
\setlength{\parindent}{2em}
\setlength{\parskip}{1.25em}
\renewcommand{\baselinestretch}{1.0}
\begin{document}
\begin{center}
{\LARGE \textit{Trump v.\ Hawaii}, 585 U.S.\ 667 (2018)\par}
\vspace{0.8em}
{\large Case Brief --- Four-Page Edition\par}
\end{center}
\vspace{1em}

\section{Memory Jogger}

Reversing a preliminary injunction against Proclamation No.~9645, the Court held 5--4 that \S~1182(f) gives the President broad power to suspend entry and that, even looking beyond the proclamation's facial neutrality under rational-basis review, the plaintiffs had not shown a likely Establishment Clause violation. It also repudiated \textit{Korematsu} (\href{https://supreme.justia.com/cases/federal/us/585/17-965/}{Justia}).

\section{Detailed Case Facts}

In 2017 President Trump issued two executive orders temporarily restricting entry from several predominantly Muslim nations and ordering a worldwide review of foreign governments' identity-management and information-sharing practices. After that review, he issued Proclamation No.~9645 under 8 U.S.C.\ \S~1182(f), which lets the President ``suspend the entry of all aliens or any class of aliens'' when he ``finds'' entry ``would be detrimental to the interests of the United States.'' The Proclamation imposed tailored restrictions on nationals of eight countries (Chad, Iran, Libya, North Korea, Somalia, Syria, Venezuela, and Yemen), exempted lawful permanent residents and asylees, offered case-by-case hardship waivers, and required reviews every 180 days. Chad was later removed, leaving five Muslim-majority nations plus North Korea and Venezuela.

\section{Procedural History}

Hawaii, three citizens with affected relatives (including Dr.\ Ismail Elshikh), and the Muslim Association of Hawaii sued, claiming the Proclamation exceeded \S~1182(f), violated the ban on nationality discrimination in immigrant-visa issuance (\S~1152(a)(1)(A)), and violated the Establishment Clause, citing the President's call for a ``total and complete shutdown of Muslims entering the United States.'' The district court entered a nationwide preliminary injunction; the Ninth Circuit affirmed on statutory grounds. The Supreme Court stayed the injunction, granted certiorari, and reversed on June 26, 2018.

\section{Judicial Votes}

Reversed and remanded, \textbf{5--4}.
\begin{itemize}
\item \textbf{Majority:} Chief Justice Roberts, joined by Justices Kennedy, Thomas, Alito, and Gorsuch.
\item \textbf{Concurrences:} Justice Kennedy; Justice Thomas.
\item \textbf{Dissents:} Justice Breyer, joined by Kagan; Justice Sotomayor, joined by Ginsburg.
\end{itemize}

\section{Holding}

The Proclamation was a lawful exercise of the President's discretion under \S~1182(f), and it did not violate the visa nondiscrimination rule of \S~1152(a)(1)(A). Under deference derived from \textit{Kleindienst v.\ Mandel} and rational-basis review, it did not violate the Establishment Clause because it was plausibly grounded in a legitimate national-security rationale independent of religious animus. \textit{Korematsu v.\ United States} was expressly repudiated.

\section{Analysis of Opinions}

\textbf{Majority (Roberts).} Assuming the statutory claims justiciable, \S~1182(f)'s only prerequisite, a finding of detriment, was met by the Proclamation's 12-page explanation of the multi-agency review, and Presidents Reagan and Clinton had imposed open-ended nationality-based suspensions. Section 1152 governs visa issuance, not the threshold question of entry, so there was no conflict. Assuming standing, the Court declined to apply domestic Establishment Clause tests: because foreign nationals abroad have no right to enter and admissions implicate national security, \textit{Mandel} asks only for a ``facially legitimate and bona fide reason.'' Even considering the President's statements, the policy survived because it could ``reasonably be understood to result from a justification independent of unconstitutional grounds'': it was neutral on its face, rested on a Cabinet-level review, had dropped three Muslim-majority countries, exempted most of the world's Muslims, and included waivers. \textit{Korematsu} ``has nothing to do with this case,'' but the Court formally repudiated it.

\textbf{Kennedy (concurring).} Officials must still honor the First Amendment even when their acts are ``not subject to judicial scrutiny or intervention.''

\textbf{Thomas (concurring).} The President has inherent power to exclude aliens, and universal injunctions exceed historical equitable power.

\textbf{Breyer (dissenting).} If the waiver system was ``window dressing,'' as tiny grant rates suggested, animus tainted the policy; he would remand while keeping the injunction.

\textbf{Sotomayor (dissenting).} A ``reasonable observer'' would see religious hostility, as in \textit{Masterpiece Cakeshop}; the majority ``redeploys the same dangerous logic'' as \textit{Korematsu}.

\textbf{Later.} \textit{Department of State v.\ Mu\~noz} (2024) relied on \textit{Hawaii} and \textit{Mandel} in rejecting a citizen's liberty interest in a spouse's visa.

\section{Examples: Future Applications}

\begin{enumerate}
\item \textbf{Same-Side: A Vetting-Based Visa Suspension.} A neutral proclamation suspending tourist visas from four countries after an interagency review, with waivers, survives despite campaign vows to ``halt Buddhist immigration.''
\item \textbf{Same-Side: Pressure Over Removals.} An open-ended suspension until a country accepts its deported nationals is lawful under \S~1182(f) and does not offend \S~1152.
\item \textbf{Opposite-Side: A Domestic Order.} Revoking accreditation eligibility for domestic colleges with Buddhist-studies programs gets full domestic Religion Clause scrutiny, without \textit{Mandel} deference.
\item \textbf{Opposite-Side: A Facially Religious Ban.} A proclamation barring adherents of a named religion cannot rest on a justification independent of unconstitutional grounds.
\item \textbf{Fence-Sitter: Paper Waivers.} Evidence that waivers were disabled and the review ignored could undermine the justification, but how much pretext overcomes deference is unsettled.
\end{enumerate}

\section{Critique}

\textbf{Double standard.} Three weeks earlier, \textit{Masterpiece Cakeshop} invalidated a decision because of commissioners' hostile remarks; here far more explicit statements were discounted. \textbf{Nondelegation.} A nearly boundless reading of \S~1182(f) sits uneasily with the Court's major-questions skepticism at home. \textbf{Symbolic overruling.} Scholars debate whether repudiating \textit{Korematsu} in dicta protects anyone if courts defer uncritically at the border.

\section{Key Quotations}

\begin{quote}
``The forcible relocation of U.S. citizens to concentration camps, solely and explicitly on the basis of race, is objectively unlawful and outside the scope of Presidential authority. ... \textit{Korematsu} was gravely wrong the day it was decided, has been overruled in the court of history, and---to be clear---`has no place in law under the Constitution.'\,''
\end{quote}
\noindent --- Chief Justice Roberts, for the Court, 585 U.S.\ at 710.

\begin{quote}
``By its terms, \S~1182(f) exudes deference to the President in every clause. It entrusts to the President the decisions whether and when to suspend entry, whose entry to suspend, for how long, and on what conditions.''
\end{quote}
\noindent --- Chief Justice Roberts, for the Court, 585 U.S.\ at 684.

\begin{quote}
``By blindly accepting the Government's misguided invitation to sanction a discriminatory policy motivated by animosity toward a disfavored group, all in the name of a superficial claim of national security, the Court redeploys the same dangerous logic underlying \textit{Korematsu} and merely replaces one `gravely wrong' decision with another.''
\end{quote}
\noindent --- Justice Sotomayor, dissenting, 585 U.S.\ at 754.

\end{document}
